% Generated draft: requires mathematical and source review.
\documentclass[11pt]{article}
\usepackage[margin=25mm]{geometry}
\usepackage{fontspec}
\setmainfont{Latin Modern Roman}
\usepackage{amsmath,amssymb,mathtools,unicode-math}
\setmathfont{Latin Modern Math}
\usepackage{enumitem,needspace,xurl,hyperref}
\setlength{\parindent}{0pt}
\setlength{\parskip}{5pt}
\setlength{\emergencystretch}{3em}
\newcommand{\N}{\mathbb{N}}
\newcommand{\Z}{\mathbb{Z}}
\newcommand{\Q}{\mathbb{Q}}
\newcommand{\R}{\mathbb{R}}
\newcommand{\C}{\mathbb{C}}
\newcommand{\F}{\mathbb{F}}
\newcommand{\A}{\mathbb{A}}
\newcommand{\PP}{\mathbb{P}}
\newcommand{\E}{\mathbb{E}}
\newcommand{\eps}{\varepsilon}
\newcommand{\dd}{\,\mathrm{d}}
\begin{document}
\section*{UnsolvedMath ID: 20000343}
\textbf{Source problem code:} catalogue entry
\subsection{Definitions and mathematical statement}
The extracted statement writes $M_{0,n}$, but the supplied research report says the official PDF has $\overline M_{0,n}$ and its overline was lost as a separate vector rule. This recovered compactification is the intended object conditionally adopted here; the open space and compactification are not interchangeable.

For $n\geq3$, $\overline M_{0,n}$ parametrizes stable ordered $n$-pointed genus-zero nodal curves, every component having at least three marked or nodal special points. Let $N_1(\overline M_{0,n})_{\R}$ be numerical curve classes, and let $\overline{NE}_1$ be the closed cone generated by effective complete curves. The original asks to determine this cone for varying $n$; the field and raw-versus-closed convention are not stated.

An $F$-curve is obtained from a partition of the markings into four nonempty sets by attaching fixed tails to a varying four-pointed rational spine, with singleton parts interpreted as markings. For part-size shape $(a,b,c,d)$, define $f_{(a,b,c,d)}$ to be the full $S_n$ orbit average, normalized as $\frac1{n!}\sum_{\sigma\in S_n}\sigma_*[F]$. This normalization matters for numerical relations.

The report treats only a symmetric slice at $n=7$. It claims
\[
 f_{(3,2,1,1)}=\tfrac12\bigl(f_{(4,1,1,1)}+f_{(2,2,2,1)}\bigr),
\]
and that $f_{(4,1,1,1)}$ and $f_{(2,2,2,1)}$ generate the two extremal rays of the invariant Mori cone. Its dual invariant nef cone is reported generated by
\[
 B_2+B_3,\qquad B_2+3B_3,
\]
where $B_i$ is understood as the sum of distinct boundary divisors whose smaller marking subset has size $i$. The exact divisor convention should be checked against the report.

It claims a complete shape/determinant audit for $n=4,5,6$ rules out earlier symmetric-shape redundancy, so among $n=4,5,6,7$ the first redundant $F$-curve shape occurs at seven. The actual enumerations and determinant matrices are not included. These unreviewed low-$n$ invariant-cone results do not determine the full nonsymmetric cone for arbitrary $n$, which the report continues to describe as open.
\subsection{Short English statement}
Determine the curve cone of the stable genus-zero moduli compactification, whose overline the extraction reportedly lost. The report computes only a symmetric seven-point cone and the first low-$n$ shape redundancy. This partial invariant calculation is distinct from the full varying-$n$ problem.
\subsection{Sources}
\begin{itemize}
\item Local UnsolvedMath registry, permanent ID $20000343$, source code \texttt{catalogue entry}; created and updated 2026-08-14. Status \texttt{partially\_solved}; classification \texttt{PARTIAL-PROGRESS}; difficulty L3; no proposer or year supplied.
\item AIM Problem Lists, Algebraic geometry, workshop \emph{Compact moduli spaces and birational geometry}, Problem 5.1, source type PDF, \url{https://aimath.org/WWN/birational/birational.pdf}. Canonical location \texttt{aim-algebraic-geometry-notes.json}, \texttt{notes[342]}, tag \texttt{problem}; statement status \texttt{exact} for the stored unbarred text. The research summary separately reports the official overline recovery. Original remarks and literature fields empty; no new source visit is claimed here.
\item Supplied research attempt 1: result \texttt{partial\_result}, run status \texttt{open}, rigor label \texttt{proved}, low novelty confidence. Candidate contribution \texttt{synthesis\_lemma}: first symmetric-shape redundancy at $n=7$, with $n=4,5,6$ enumeration and determinant certificates. No exact primary references or audit matrices are supplied. Explicit caveat: machine-generated research aid, not peer reviewed; independently verify before citation.
\item Dataset attribution: Ulam AI / UnsolvedMath Contributors, CC BY 4.0; original problem attribution remains to the supplied AIM workshop collection.
\end{itemize}
\end{document}
